\section{Introduzione e obiettivi}
\label{sec:introduction}

\subsection{Contesto e obiettivi del progetto}

Il progetto realizza un prototipo distribuito per monitorare il
livello dell'acqua piovana in una tanica e controllare l'apertura
di un canale di scarico. Il sistema prevede l'acquisizione periodica
delle misure, il controllo automatico dell'apertura e la possibilità
di intervento da parte dell'operatore, sia localmente sia attraverso
un'interfaccia web.

La consegna dell'Assignment 03 definisce i requisiti funzionali,
i sottosistemi e le modalità di comunicazione previste~\cite{assignment03}.
Il sistema è composto da quattro sottosistemi:
il Tank Monitoring Subsystem (TMS), responsabile dell'acquisizione
delle misure; il Control Unit Subsystem (CUS), che coordina il
funzionamento del sistema; il Water Channel Subsystem (WCS), che
gestisce l'attuatore e l'interfaccia locale; e il Dashboard
Subsystem (DBS), che permette il monitoraggio e il controllo remoto.

\subsection{Comportamento e requisiti principali}

Il sistema si avvia in modalità \code{AUTOMATIC}, nella quale
l'apertura del canale viene regolata in base al livello dell'acqua
e alle soglie temporali previste. In modalità \code{MANUAL},
l'operatore può impostare direttamente l'apertura desiderata.
È inoltre prevista la gestione dell'indisponibilità delle misure
e delle comunicazioni.

Il prototipo utilizza un ESP32 per il TMS, un Arduino UNO per
il WCS, un backend Java eseguito su PC per il CUS e un'interfaccia
web per il DBS. Le comunicazioni con il CUS avvengono rispettivamente
tramite MQTT, connessione seriale e HTTP.
